\section{배경}
\label{sec:background}

대부분의 PUBG 예측 연구는 경기 종료 후 통계를 이용해 최종 순위를
예측한다~\cite{kagglePubgFinishPlacement,sen2021pubgFinalRank,lee2025pubgRank}.
이 접근은 최종 결과 모델링에는 적합하지만, 위험이 발생한 시간적
순서와 공간 맥락을 제거한다. 가장 가까운 선행 연구는 Sagong 등의
페이즈 기반 프레임워크이다. 이 연구는 FightScore와 ActiveScore로
플레이어의 공격성과 활동성을 점수화하고, 전략 유형을 군집화한 뒤
다음 페이즈 생존을 예측한다~\cite{sagong2025hunterCamper}. 우리는
페이즈 수준 문제 설정을 계승하되, 분석 단위를 플레이어에서 팀으로
바꾸고 작은 전략 점수 요약 대신 전체 텔레메트리 재구성을 사용한다.

수축하는 안전 구역은 시공간적 관점도 요구한다. 서식지 단편화와
movement ecology는 이용 가능한 공간, 이동, 노출, 밀도를 서로 결합된
시간-공간 변수로 다룬다
~\cite{fahrig2003fragmentation,nathan2008movement}. 우리는 이를
생물학적 주장으로 사용하지 않고 분석적 비유로만 사용한다. PUBG 팀은
유기체가 아니지만, 그들의 생존 조건은 축소되는 활동 공간에 의해
형성된다. 따라서 각 페이즈를 특징 행과 그래프 스냅샷으로 동시에
표현한다. 그래프 edge는 어느 팀이 가까운지, 어느 팀이 누구를
공격했는지, 어떤 팀들이 같은 안전존 압력을 받는지를 인코딩한다
~\cite{sanchezLengeling2021gnnIntro}.
